\chapter{Reproducibility}
\label{app:reproducibility}

\section{Abstract}

The artefacts supporting this thesis are distributed across three public Git repositories: an infrastructure repository that provisions a single-node Kubernetes 1.34 cluster on Hetzner Cloud via Terraform plus \texttt{cloud-init} plus \texttt{kubeadm}; an experiments repository containing the four benchmark variants (\texttt{wasm-rust}, \texttt{wasm-tinygo}, \texttt{docker-rust}, \texttt{docker-golang}) across the two stateless \acs{HTTP} workloads (prime sieve, memory bandwidth), the k6 load-generator scripts, and the Prometheus / chart-regeneration helpers; and this thesis source. The end-to-end reproduction path is two \texttt{make} invocations against the infrastructure repository followed by a per-benchmark \texttt{run\_experiment.sh} script in the experiments repository, on a single Hetzner ccx23 \acs{VM} (4 vCPU, 16\,GB \acs{RAM}). The multi-node scaling experiments referenced as future work in Section~\ref{sec:future_work_multinode} are deferred and are not part of this artefact.

\section{Artifact check-list (meta-information)}

\begin{description}
    \item[Program:] Sieve of Eratosthenes \acs{HTTP} service (benchmark 01) and heap-allocation~+~SHA-256 \acs{HTTP} service (benchmark 02), each implemented as four variants (\texttt{wasm-rust}, \texttt{wasm-tinygo}, \texttt{docker-rust}, \texttt{docker-golang}).
    \item[Compilation:] Rust 1.94 with \texttt{cargo} and \texttt{cargo-component}, Go 1.26, TinyGo 0.40.0 (\texttt{-target=wasip1 -gc=conservative -opt=2}).
    \item[Run-time environment:] Ubuntu 24.04 LTS, Kubernetes 1.34 via \texttt{kubeadm}, containerd 2.0.2, \texttt{containerd-shim-spin-v2} v0.17.0, Flannel CNI, local-path-provisioner.
    \item[Hardware:] One Hetzner Cloud ccx23 \acs{VM} (4$\times$ AMD EPYC vCPU, 16\,GB \acs{RAM}, 160\,GB NVMe), location \texttt{nbg1}.
    \item[Metrics:] Requests per second; end-to-end latency p50/p95/p99; cold- and warm-start latency; container \acs{RSS} at idle and at peak load; \acs{CPU} consumption in millicores; on-disk \acs{OCI} artifact size.
    \item[Output:] Per-variant \texttt{*\_summary.json} files (aggregate k6 metrics), per-variant \texttt{*\_k6.json} files (raw k6 output), and per-mode aggregates \texttt{cold\_start.json}, \texttt{warm\_start.json}, \texttt{resource\_metrics.json}, \texttt{image\_sizes.json}, together with the PNG charts under \texttt{thesis-experiments/results/0\{1,2\}-*/\{limited,unlimited\}/} that Chapter~\ref{chap:evaluation} embeds.
    \item[Experiments:] \texttt{01-prime-sieve} and \texttt{02-memory-bandwidth}, each run in both \emph{limited} (single worker thread / single replica) and \emph{unlimited} (\texttt{GOMAXPROCS=4}, \texttt{TOKIO\_WORKER\_THREADS=4}, \texttt{SpinApp~spec.replicas:~4}) modes.
    \item[Approximate disk space:] A few \acs{GB} total, dominated by raw k6 \acs{JSON} output and the Prometheus \acs{PVC}.
    \item[Approximate wall-clock to reproduce:] ${\sim}30$\,min for cluster provisioning, plus ${\sim}20$\,min per benchmark mode, totalling ${\sim}2$\,h end-to-end.
    \item[Publicly available:] Yes; the three repositories are cited as \cite{musaev2026thesisinfra, musaev2026thesisexperiments, musaev2026thesissource}.
\end{description}

\section{Description}
\subsection{How to access}
\label{sec:repro_how_to_access}

The complete set of artefacts supporting the experiments in this thesis is distributed across three public Git repositories hosted on GitHub. Each repository is referenced from its corresponding chapter; this subsection groups them together as a single point of entry.

\begin{itemize}
    \item \textbf{Infrastructure setup} \textendash{} \cite{musaev2026thesisinfra}.
          Kubernetes (\texttt{kubeadm}) and SpinKube cluster provisioning, Hetzner Cloud Terraform modules, cloud-init bootstrap scripts, and the \texttt{containerd}~/ \texttt{containerd-shim-spin-v2} installation manifests used to build the four-variant experimental cluster described in \hyperref[chap:implementation]{Chapter~\ref{chap:implementation}}.

    \item \textbf{Benchmark implementations and result data} \textendash{} \cite{musaev2026thesisexperiments}.
          The four primary benchmark variants (\texttt{docker-golang}, \texttt{docker-rust}, \texttt{wasm-rust}, \texttt{wasm-tinygo}), the k6 load-generator configurations, the Prometheus scrape rules, and the raw \texttt{results/} directory referenced throughout \hyperref[chap:evaluation]{Chapter~\ref{chap:evaluation}}.

    \item \textbf{Thesis source} \textendash{} \cite{musaev2026thesissource}.
          The \LaTeX{} source of this document, including all chapter \texttt{.tex} files, the bibliography, TikZ figure sources, and the build instructions in the \texttt{Makefile}.
\end{itemize}

All three repositories are tagged at the commit corresponding to thesis submission so that a reader can clone the exact state of the artefacts that produced the reported results.

\subsection{Hardware dependencies}
\label{sec:repro_hardware}

The single-node measurement environment requires one Hetzner Cloud ccx23 \acs{VM} (4$\times$ AMD EPYC vCPU, 16\,GB \acs{RAM}, 160\,GB NVMe \acs{SSD}). The Hetzner default 8-core account quota is sufficient. The Terraform module accepts the location \texttt{nbg1} (Nuremberg) by default; any other Hetzner Cloud location should work without changes. An admin workstation with outbound \acs{SSH} (port 22) and \texttt{kubectl} (port 6443) connectivity to the \acs{VM}'s public \acs{IP} is required for provisioning and benchmark execution. The deferred multi-node mode (Section~\ref{sec:future_work_multinode}) would require additional Hetzner instances under the same account, but is not exercised by this artefact.

\subsection{Software dependencies}
\label{sec:repro_software}

On the admin workstation: Terraform $\geq 1.6$, \texttt{kubectl} matching cluster version 1.34, Helm~3, \texttt{make}, Docker~26 (to build and push the two Docker variants to a registry), Rust~1.94 with \texttt{rustup target add wasm32-wasip2} and \texttt{cargo install cargo-component --locked}, Go~1.26, TinyGo~0.40.0 (installed from the upstream \texttt{.tar.gz} release; no stable \texttt{apt} package on Ubuntu), the Spin \acs{CLI} v3+, k6 (installable via the Grafana \texttt{apt} repository), and Python~3 with \texttt{matplotlib} for chart regeneration.

On the provisioned cluster, the following components are installed automatically by \texttt{cloud-init} and by the \texttt{make deploy} target and therefore do not need to be installed manually: containerd~2.0.2, \texttt{runc}~1.2.4, CNI plugins~1.6.2, \texttt{containerd-shim-spin-v2}~v0.17.0, \texttt{cert-manager}~v1.16.3, SpinOperator~v0.6.1 (Helm chart \texttt{spinframework/spin-operator}), and \texttt{kube-prometheus-stack} (community Helm chart, with the values in \texttt{observability-values.yaml}).

\section{Installation}
\label{sec:repro_installation}

The following procedure provisions the single-node cluster end-to-end. All \texttt{make} targets referenced below are defined in \texttt{thesis-infra-setup/Makefile} and are idempotent: re-running \texttt{make deploy} against an established cluster is a no-op.

\begin{enumerate}
    \item Clone the three repositories side-by-side under a common parent directory (for example \texttt{\textasciitilde/coding/master-thesis/}). The Terraform Makefile and the per-benchmark \texttt{run\_experiment.sh} scripts assume this sibling layout.
    \item Populate \texttt{thesis-infra-setup/terraform.tfvars} with a Hetzner Cloud API token (\texttt{hcloud\_token}), an \acs{SSH} public key (\texttt{ssh\_public\_key}), and the admin workstation's egress \acs{IP}/CIDR (\texttt{admin\_ip\_cidr}). Leave \texttt{worker\_count = 0}, which is the single-node default.
    \item Provision the cluster: \texttt{make up}. This runs \texttt{terraform apply}, waits for \texttt{cloud-init} to complete, runs \texttt{kubeadm init} for Kubernetes 1.34, installs the Flannel CNI plug-in and the local-path-provisioner, registers \texttt{containerd-shim-spin-v2} with containerd, and removes the control-plane taint so that the single node can also schedule workload pods.
    \item Fetch the cluster's kubeconfig: \texttt{make configure}. This writes \texttt{hetzner-thesis.yaml} into the repository root and rewrites the API-server endpoint to the \acs{VM}'s public \acs{IP}.
    \item Label the node for SpinKube pod scheduling: \texttt{make label}.
    \item Install the cluster add-ons: \texttt{make deploy}. This installs \texttt{cert-manager}~v1.16.3, the SpinOperator Helm chart (v0.6.1) together with the \texttt{wasmtime-spin} \texttt{RuntimeClass}, and \texttt{kube-prometheus-stack} parameterised by \texttt{observability-values.yaml} (5\,s scrape interval, Grafana exposed as a \acs{NodePort}).
    \item Sanity-check the runtime path: \texttt{make test} applies \texttt{test-spin.yaml} (a Spin ``hello'' sample) and expects an \acs{HTTP}~200 response.
\end{enumerate}

Steps~1--7 typically complete in ${\sim}30$\,minutes of wall-clock time.

\section{Experiment workflow}
\label{sec:repro_workflow}

The two benchmarks share the same four-variant structure, the same NodePort assignment (30081 \texttt{wasm-rust}, 30082 \texttt{wasm-tinygo}, 30083 \texttt{docker-rust}, 30084 \texttt{docker-golang}), and the same orchestrator script layout. They are run sequentially because the two benchmarks share NodePorts; the \texttt{02-memory-bandwidth} script automatically tears down the \texttt{01-prime-sieve} \texttt{SpinApp}s and \texttt{Deployment}s before deploying its own.

\subsection{Benchmark 01 --- prime sieve}
\label{sec:repro_workflow_01}

Build and push the four variant images:

\begin{itemize}
    \item \textbf{\texttt{wasm-rust}:} \texttt{cd wasm/rust/01-prime-sieve \&\& cargo component build --release} followed by \texttt{spin registry push docker.io/<user>/prime-sieve-wasm-rust:latest}. Build path: \texttt{thesis-experiments/wasm/rust/01-prime-sieve/}.
    \item \textbf{\texttt{wasm-tinygo}:} \texttt{cd wasm/tinygo/01-prime-sieve \&\& tinygo build -target=wasip1 -gc=conservative -opt=2 -o app.wasm .} followed by \texttt{spin registry push docker.io/<user>/prime-sieve-wasm-tinygo:latest}. Build path: \texttt{thesis-experiments/wasm/tinygo/01-prime-sieve/}.
    \item \textbf{\texttt{docker-rust}:} \texttt{docker build -t <reg>/prime-sieve-docker-rust:latest docker/rust/01-prime-sieve/} followed by \texttt{docker push}. The \texttt{Dockerfile} uses \texttt{rust:1.94-*} as the builder and a \texttt{scratch} runtime stage.
    \item \textbf{\texttt{docker-golang}:} \texttt{docker build -t <reg>/prime-sieve-docker-golang:latest docker/golang/01-prime-sieve/} followed by \texttt{docker push}. The \texttt{Dockerfile} uses \texttt{golang:1.26-alpine} as the builder.
\end{itemize}

Deploy the four variants: \texttt{kubectl apply -f k8s/01-prime-sieve/namespace.yaml \&\& kubectl apply -f k8s/01-prime-sieve/}. The manifests bind the four NodePorts listed above; the variant order matches the NodePort-ascending convention used throughout this thesis.

Run the experiment: \texttt{./benchmarks/01-prime-sieve/run\_experiment.sh -{}-scaling-experiment both}. The script exports \texttt{KUBECONFIG=../thesis-infra-setup/hetzner-thesis.yaml}, reads the node's public \acs{IP} from \texttt{terraform output -raw instance\_public\_ip}, sequentially drives k6 (50 \acp{VU}, 20\,s ramp, 60\,s hold) against each variant in both the \emph{limited} and \emph{unlimited} configurations, queries Prometheus for \texttt{container\_memory\_rss} and \texttt{container\_cpu\_usage\_seconds\_total}, runs the cold-start (\texttt{kubectl scale --replicas=0} $\to$ 1) and warm-start measurement sequences, captures \acs{OCI} artifact sizes via \texttt{docker inspect} (and, for the Spin variants, via the compiled \texttt{.wasm} artifact on disk), and finally invokes \texttt{analyze.py} to regenerate the per-mode PNG charts under \texttt{results/01-prime-sieve/\{limited,unlimited\}/}.

\subsection{Benchmark 02 --- memory bandwidth}
\label{sec:repro_workflow_02}

The procedure is identical in shape. Build paths are \texttt{thesis-experiments/\{docker,wasm\}/\{rust,golang,tinygo\}/02-memory-bandwidth/}. Deploy with \texttt{kubectl apply -f k8s/02-memory-bandwidth/}; the script first tears down any \texttt{01-prime-sieve} pods because both benchmarks share the four NodePorts. Run with \texttt{./benchmarks/02-memory-bandwidth/run\_experiment.sh -{}-scaling-experiment both}. The endpoint is \texttt{GET /membw?size\_kb=64\&no\_hash=0}; the k6 ramp / hold parameters, the Prometheus queries, and the cold/warm-start sequence mirror benchmark~01.

The limited-mode configuration uses \texttt{GOMAXPROCS=1} (\texttt{docker-golang}), \texttt{TOKIO\_WORKER\_THREADS=1} (\texttt{docker-rust}), and \texttt{SpinApp~spec.replicas:~1} (both \acs{Wasm} variants). The unlimited-mode configuration raises these to \texttt{GOMAXPROCS=4}, \texttt{TOKIO\_WORKER\_THREADS=4}, and \texttt{SpinApp~spec.replicas:~4}, and additionally raises the Kubernetes \texttt{limits.cpu} to \texttt{4000m} so that cgroup CPU bandwidth control does not throttle the additional worker threads in the 4-vCPU environment of the Hetzner ccx23 host.

\section{Evaluation and expected results}
\label{sec:repro_evaluation}

Each experiment run materialises a deterministic on-disk layout under \texttt{thesis-experiments/results/0\{1,2\}-*/\{limited,unlimited\}/}: per-variant \texttt{*\_summary.json} files containing the aggregated k6 metrics, per-variant \texttt{*\_k6.json} files containing the raw k6 output for percentile recomputation, and the four per-mode aggregates \texttt{cold\_start.json}, \texttt{warm\_start.json}, \texttt{resource\_metrics.json}, and \texttt{image\_sizes.json}. The PNG charts that \texttt{chapters/06\_evaluation.tex} embeds are written into the same directory by \texttt{analyze.py}.

A successful reproduction should yield, on \texttt{01-prime-sieve} in limited mode, approximately \texttt{docker-golang}~3250\,\acs{RPS} (p50 11.6\,ms), \texttt{wasm-rust}~3058\,\acs{RPS} (p50 12.8\,ms), \texttt{docker-rust}~2073\,\acs{RPS} (p50 20.7\,ms), and \texttt{wasm-tinygo}~1419\,\acs{RPS} (p50 29.5\,ms); on \texttt{02-memory-bandwidth} in limited mode, the four variants fall within the brackets reported in Section~\ref{sec:eval_performance}. Cold-start latency for both \acs{Wasm} variants and \texttt{docker-rust} sits in a ${\sim}25$--$35$\,ms band; \texttt{docker-golang}'s cold-start measurement should land in the multi-second range (${\sim}3.7$--$4.8$\,s). Absolute \acs{RPS} values may shift by a few percent between runs because of per-host EPYC-core variance, but the relative ordering of the four variants and the order-of-magnitude artifact-size gap (\texttt{wasm-rust} ${\sim}30\times$ smaller than \texttt{docker-golang}) are stable.

The chapter-6 figures can be regenerated \emph{without} re-running the benchmarks by invoking \texttt{python3 benchmarks/01-prime-sieve/analyze.py -{}-mode limited} and \texttt{-{}-mode unlimited}, and the equivalents under \texttt{benchmarks/02-memory-bandwidth/}. The script reads the existing \acs{JSON} files under \texttt{results/} and emits the exact PNG filenames that the evaluation chapter includes.

\section{Notes}
\label{sec:repro_notes}

The single-node measurement environment was chosen so that the four-variant comparison isolates runtime overhead from inter-node scheduling, kube-proxy load balancing, and network-fabric latency. The chapter-6 numbers therefore characterise per-pod behaviour rather than cross-host distributed behaviour; the deferred multi-node validation (different host fabric, replicas placed on distinct workers) is scoped in Section~\ref{sec:future_work_multinode} and is intentionally not part of this artefact.

Wasmtime runs in its default \acf{JIT} mode (Cranelift) throughout. Wasmtime's \acf{AOT} pre-compilation pathway via the \texttt{wasmtime compile} command is expected to narrow the throughput gap with native code, as discussed in Section~\ref{sec:future_work_aot}; the reported numbers therefore represent a conservative lower bound on the Wasmtime throughput attainable on this hardware.

The four variants' source code, \texttt{Dockerfile}s, \texttt{SpinApp} \acs{CRD} manifests, and \texttt{spin.toml} configurations are also discussed in \hyperref[chap:implementation]{Chapter~\ref{chap:implementation}}. This appendix focuses on the operational view (commands, version pins, output shape, regeneration path) so that the chapter prose and the appendix do not duplicate hardware and software detail.
